\documentclass[10pt,a4paper]{amsart}
\usepackage[margin=19mm]{geometry}
\usepackage{amsmath,amssymb,mathtools}
\usepackage[scr=rsfs,cal=euler]{mathalfa}
\usepackage{xfrac}
\usepackage[unicode,hidelinks,bookmarks=false]{hyperref}
\newcommand{\ie}{i.e.\ }
\newcommand{\cf}{cf.\ }
\newcommand{\resp}{respectively\ }
\newcommand{\Subsection}{\subsection}
\providecommand{\nobreakdash}{\nobreak\hbox{-}}
\setlength{\emergencystretch}{2em}
\multlinegap0pt
\usepackage{amssymb}
\newtheorem{definition}{Definition}
\newtheorem{lemma}{Lemma}
\title{On a local variant of the 12th Delfino problem --- the $\Pi$-side}
\author{Stefan Hoffelner}
\date{Source: arXiv:2606.13210v1 (CC BY 4.0)}
\begin{document}
\maketitle

\noindent\emph{Source and licence.} On a local variant of the 12th Delfino problem --- the $\Pi$-side. Version arXiv:2606.13210v1 (CC BY 4.0), distributed by arXiv under the Creative Commons Attribution 4.0 International licence, http://creativecommons.org/licenses/by/4.0/, as recorded in the arXiv licence metadata for this exact version and shown on the abstract page https://arxiv.org/abs/2606.13210v1. Full licence notice: \texttt{licenses/arxiv-2606.13210-CC-BY-4.0.txt}.\par\smallskip

\noindent\textit{Editorial scope.} Counted unit: the article's lemma lem:rho-allowable-products-local (source0.tex lines 2851--2874). The article's complete original proof of it is delivered with it (source0.tex lines 2876--3034, one \texttt{proof} environment of 159 lines). No mathematics is written, repaired or re-derived: the statement and the proof are the article's own lines, in the article's own notation, and no retained line is substituted or re-flowed.  Retained uncounted as the source-local context the proof needs: lines 941--957; lines 1638--1732; lines 1839--2012; lines 2626--2734; lines 2736--2746; lines 2799--2820.  Nothing was omitted from any retained span, so the package carries no omission record.  No retained line contains a citation, so the wrapper carries no bibliography and no bibliography entry.  No retained line contains a figure environment, an image-inclusion command or drawing code, so no image is delivered and the package declares no asset dependency.  Editorial formatting only, in addition: an AMS \texttt{amsart} preamble that restates the article's own declarations and macros used by the retained lines, namely the environments \texttt{definition}, \texttt{lemma} on separate counters, together with the standard packages those lines need.  The retained environments print in this document as Definition 1, Lemma 1 in order of appearance, whereas the article prints the same items with the numbering of its own full document.  The complete native source of the article as distributed in the arXiv e-print for this exact version is bundled unchanged as \texttt{sources/202684/source0.tex}.
\begin{definition}[Selected branch coordinates]\label{def:selected-branch-coordinates}
For \(C\subseteq\omega_1\) and \(u\in2^\omega\), define
\[
   \operatorname{Sel}(C,u)=
   \{\tau_u(\gamma,n)\mid \gamma\in C,\ n<\omega\},
\]
where
\[
   \tau_u(\gamma,n)=
   \begin{cases}
      \omega\cdot\gamma+2n, &\text{if } n\notin u,\\
      \omega\cdot\gamma+2n+1, &\text{if } n\in u.
   \end{cases}
\]
Thus \(\operatorname{Sel}(C,u)\) is the set of branch coordinates selected by
the area \(C\) when it writes the real \(u\).
\end{definition}

\begin{definition}[Regularity-stage supports and branch footprints]
\label{def:admissible-random-base}
Let \(\mathcal R_\beta\) be a local hybrid presentation over \(W=M_1[H]\), and
let \(\eta<\beta\) be a relative regularity stage.  Suppose that the resolved
tag at \(\eta\) is one of
\[
   (\mathrm{Rand};\dot a^\eta_0,\ldots,\dot a^\eta_{k_\eta-1}),
   \qquad
   (\mathrm{Am}_{\mathcal N};\dot a^\eta_0,\ldots,
      \dot a^\eta_{k_\eta-1}),
   \qquad
   (\mathrm{Am}_{\mathcal M};\dot a^\eta_0,\ldots,
      \dot a^\eta_{k_\eta-1}).
\]
The base used at this stage is, by definition,
\[
   \dot W_\eta=L[T_2,\dot a^\eta_0,\ldots,
      \dot a^\eta_{k_\eta-1}],
\]
with the evident meaning \(L[T_2]\) when \(k_\eta=0\).  The regularity forcing
at the stage is the Random algebra, measure-amoeba forcing, or category-amoeba
forcing computed in this displayed base, according to the tag.

A \emph{support witness} for the regularity stage is a tuple
\[
   (A_{\mathrm{rb}},
    \langle \dot a^{\mathrm{rb}}_i\mid i<k_\eta\rangle,
    B^0_{\mathrm{rb}},
    \dot B_{\mathrm{rb}})
\]
such that:
\begin{enumerate}
   \item \(A_{\mathrm{rb}}\subseteq\eta\) is a countable admissible support, and
   each \(\dot a^{\mathrm{rb}}_i\) is a \(\mathbb P_{A_{\mathrm{rb}}}\)-name for
   a real;

   \item for every \(i<k_\eta\),
   \[
      \mathbb P_\eta\Vdash
      \dot a^\eta_i=(\dot a^{\mathrm{rb}}_i)^{\uparrow\eta};
   \]

   \item every object of \(W=M_1[H]\) occurring in the codes of the names
   \(\dot a^{\mathrm{rb}}_i\) is represented by its \(<_{M_1}\)-least
   \(Br\)-name, and the union of the supports of these \(Br\)-names is
   contained in the countable set \(B^0_{\mathrm{rb}}\subseteq\omega_1\).
\end{enumerate}
Let
\[
   \dot W_{\mathrm{rb}}
   =L[T_2,
      \dot a^{\mathrm{rb}}_0,\ldots,
      \dot a^{\mathrm{rb}}_{k_\eta-1}]
\]
be the base read on the support \(A_{\mathrm{rb}}\).  Thus the canonical lift of
\(\dot W_{\mathrm{rb}}\) to stage \(\eta\) is forced to be the stage base
\(\dot W_\eta\).

Using the notation from Definition~\ref{def:selected-branch-coordinates}, let
\(E(A_{\mathrm{rb}})\) be the set of explicit coding coordinates in
\(A_{\mathrm{rb}}\).  If \(\nu\in E(A_{\mathrm{rb}})\), let \(\dot C_\nu\) be the
\(\mathbb P_{A_{\mathrm{rb}}}\)-name for the coding area of the reservoir
generic at \(\nu\), and let \(\dot u_\nu\) be the
\(\mathbb P_{A_{\mathrm{rb}}}\)-name coded at \(\nu\).  The branch footprint of
the witness is the \(\mathbb P_{A_{\mathrm{rb}}}\)-name
\[
   \dot B_{\mathrm{rb}}
   =
   \check B^0_{\mathrm{rb}}
   \cup
   \bigcup_{\nu\in E(A_{\mathrm{rb}})}
      \operatorname{Sel}(\dot C_\nu,\dot u_\nu).
\]
Thus \(\dot B_{\mathrm{rb}}\) is generated by countably many coding areas and a
countable set of branch coordinates.  It may have size \(\omega_1\).

Among all support witnesses for the stage, choose the \(<_{M_1}\)-least one.  We
call the associated quadruple
\[
   (A^{\mathrm{reg}}_\eta,
    \dot W^{\mathrm{reg}}_\eta,
    B^{0,\mathrm{reg}}_\eta,
    \dot B^{\mathrm{reg}}_\eta)
\]
the \emph{canonical support witness} for the regularity stage, where
\[
   \dot W^{\mathrm{reg}}_\eta
   =L[T_2,
      \dot a^{\mathrm{reg}}_0,\ldots,
      \dot a^{\mathrm{reg}}_{k_\eta-1}]
\]
is the base read on the least support witness.  The regularity forcing is always
computed in the relevant \(L[T_2,\vec a]\)-base.  It is not recomputed after
passing to a larger ambient extension.
\end{definition}

\begin{definition}[\(0\)-allowable local hybrid presentations]
\label{def:zero-local-allowable}
Let \(F\in M_1\) be a bookkeeping function of length
\(\delta\leq\omega_2\).  A sequence
\[
   \mathcal R=\langle \mathcal R_\beta\mid \beta\leq\delta\rangle
\]
over \(W\) is a \emph{\(0\)-allowable local hybrid presentation relative to
\(F\)} if it is obtained by the following recursion.

First, \(\mathcal R_0\) is the trivial post-branch presentation.  Thus the associated quotient forcing over \(W\) is
\[
   \mathbb P_0=\mathbf 1,
\]
whereas the corresponding full forcing over \(M_1\) is simply \(Br\).
At limit stages, \(\mathcal R_\eta\) is the mixed-support limit of the preceding
hybrid construction, and \(\mathbb P_\eta\) is the associated post-branch
forcing.  Suppose that \(\mathcal R_\beta\) has been constructed, and write
\(\mathbb P_\beta\) for its associated post-branch forcing.
Choose a fresh reservoir coordinate \(\mathbb C_\beta\), let
\(\dot g_\beta\) be its canonical generic, and put
\[
   \widehat{\mathbb P}_{\beta+1}
      =
   \mathbb P_\beta\times\mathbb C_\beta.
\]
The second-step forcing \(\dot{\mathbb Q}_\beta\) over
\(\widehat{\mathbb P}_{\beta+1}\) is determined by the decoded value of
\(F(\beta)\).

\begin{enumerate}
   \item If \(F(\beta)\) is one of the relative regularity tags
   \[
      (\mathrm{Rand};\dot a^\beta_0,\ldots,
         \dot a^\beta_{k_\beta-1}),
      \qquad
      (\mathrm{Am}_{\mathcal N};\dot a^\beta_0,\ldots,
         \dot a^\beta_{k_\beta-1}),
      \qquad
      (\mathrm{Am}_{\mathcal M};\dot a^\beta_0,\ldots,
         \dot a^\beta_{k_\beta-1}),
   \]
   and this stage has a canonical admissibility witness in the sense of
   Definition~\ref{def:admissible-random-base}, let
   \[
      (A^{\mathrm{reg}}_\beta,
       \dot W^{\mathrm{reg}}_\beta,
       B^{0,\mathrm{reg}}_\beta,
       \dot B^{\mathrm{reg}}_\beta)
   \]
   be that witness.  Thus, for the \(<_{M_1}\)-least tuple of
   \(\mathbb P_{A^{\mathrm{reg}}_\beta}\)-names for reals
   \(\langle\dot a^{\beta,\mathrm{reg}}_i\mid i<k_\beta\rangle\) reading the
   displayed real names, we have
   \[
      \dot W^{\mathrm{reg}}_\beta
      =L[T_2,
          \dot a^{\beta,\mathrm{reg}}_0,
          \ldots,
          \dot a^{\beta,\mathrm{reg}}_{k_\beta-1}].
   \]
   Let \(\dot W^{\mathrm{reg},+}_\beta\) be the canonical
   \(\widehat{\mathbb P}_{\beta+1}\)-name obtained by lifting this base, equivalently
   the name for
   \[
      L[T_2,
        (\dot a^{\beta,\mathrm{reg}}_0)^{\uparrow},\ldots,
        (\dot a^{\beta,\mathrm{reg}}_{k_\beta-1})^{\uparrow}]
   \]
   over the stage \(\widehat{\mathbb P}_{\beta+1}\).  The second-step forcing
   \(\dot{\mathbb Q}_\beta\) is the member of the following list corresponding
   to the tag:
   \[
      \dot{\mathbb B}_{\mathrm{rand}}^{\dot W^{\mathrm{reg},+}_\beta},
      \qquad
      \dot{\mathbb A}_{\mathcal N}^{\dot W^{\mathrm{reg},+}_\beta},
      \qquad
      \dot{\mathbb A}_{\mathcal M}^{\dot W^{\mathrm{reg},+}_\beta}.
   \]
   Thus the relative Random or amoeba forcing is computed in the lifted base
   \(L[T_2,(\dot a^{\beta,\mathrm{reg}}_i)^{\uparrow}:i<k_\beta]\).  It is not
   recomputed in the larger ambient extension.  The local support
   \(A^{\mathrm{reg}}_\beta\) and the footprint name
   \(\dot B^{\mathrm{reg}}_\beta\) are recorded as data of this stage.  If the
   displayed regularity tag has no admissibility witness, put
   \(\dot{\mathbb Q}_\beta=\mathbf 1\).

   \item If \(F(\beta)\) is an ordinary Cohen tag, then
   \[
      \dot{\mathbb Q}_\beta=\dot{\mathbb C}_\omega.
   \]

   \item If \(F(\beta)=\dot w\), where \(\dot w\) is a
   \(\mathbb P_\beta\)-name for a real, let
   \[
      \dot u_\beta=\dot w
   \]
   and put
   \[
      \dot{\mathbb Q}_\beta
      =
      \dot{\mathbb A}_D(\dot Y_{\dot g_\beta,\dot u_\beta}).
   \]

   \item If
   \[
      F(\beta)=(\dot x,\dot z,\dot m),
   \]
   where \(\dot x,\dot z\) are \(\mathbb P_\beta\)-names for reals and
   \(\dot m\) is a \(\mathbb P_\beta\)-name for a natural number, let
   \[
      \dot u_\beta=\langle \dot x,\dot z,\dot m\rangle
   \]
   and put
   \[
      \dot{\mathbb Q}_\beta
      =
      \dot{\mathbb A}_D(\dot Y_{\dot g_\beta,\dot u_\beta}).
   \]

   \item If
   \[
      F(\beta)=(\dot z_0,\dot z_1),
   \]
   where \(\dot z_0,\dot z_1\) are \(\mathbb P_\beta\)-names for reals, let
   \((A_i,\sigma_i)\) be the canonical localized presentation of \(\dot z_i\),
   for \(i<2\).  If
   \[
      (A_0,\sigma_0)<_{M_1}(A_1,\sigma_1),
   \]
   let
   \[
      \dot u_\beta=\langle \dot z_0,\dot z_1\rangle
   \]
   and put
   \[
      \dot{\mathbb Q}_\beta
      =
      \dot{\mathbb A}_D(\dot Y_{\dot g_\beta,\dot u_\beta}).
   \]
   If instead
   \[
      (A_1,\sigma_1)<_{M_1}(A_0,\sigma_0),
   \]
   let
   \[
      \dot u_\beta=\langle \dot z_1,\dot z_0\rangle
   \]
   and put
   \[
      \dot{\mathbb Q}_\beta
      =
      \dot{\mathbb A}_D(\dot Y_{\dot g_\beta,\dot u_\beta}).
   \]
   If the two canonical localized presentations are equal, put
   \[
      \dot{\mathbb Q}_\beta=\mathbf 1.
   \]
\end{enumerate}
If none of the preceding cases applies, put
\[
   \dot{\mathbb Q}_\beta=\mathbf 1.
\]
Finally define
\[
   \mathbb P_{\beta+1}
      =
   \widehat{\mathbb P}_{\beta+1}*\dot{\mathbb Q}_\beta.
\]
In cases (3), (4), and the nontrivial subcases of (5), we say that stage
\(\beta\) \emph{explicitly codes} the name \(\dot u_\beta\).  A forcing
admitting such a presentation is called \emph{locally allowable}, or
\emph{\(0\)-allowable}.
\end{definition}

\begin{definition}[Successor \(\rho\)-allowability]
\label{def:rho-allowable-local}
Assume that \(\rho=\alpha+1\) and that \(\zeta\)-allowability and relative
\(\zeta\)-allowability have been defined for every \(\zeta<\rho\).  Let
\(F\in M_1\) be a bookkeeping function of length \(\delta\leq\omega_2\).  We
say that \((\mathcal R_\delta,\dot I_\delta)\) is \emph{\(\rho\)-allowable
relative to \(F\)} if it is obtained by the following recursion in \(W\).

At each stage \(\beta\leq\delta\) we construct a hybrid presentation
\(\mathcal R_\beta\) and a \(\mathbb P_\beta\)-name \(\dot I_\beta\), where
\(\mathbb P_\beta\) denotes the post-branch forcing associated with
\(\mathcal R_\beta\).
We start with $W$ and with \(\dot I_0\) the
canonical name for the empty set.  At a limit stage \(\eta\leq\delta\),
\(\mathcal R_\eta\) is the mixed-support limit of the preceding hybrid
presentations, and \(\dot I_\eta\) is the canonical \(\mathbb P_\eta\)-name
forced to be the union of the earlier interpreted auxiliary names.

Suppose that \(\beta<\delta\) and that \((\mathcal R_\beta,\dot I_\beta)\) has
already been constructed.  The value \(F(\beta)\) is decoded relative to
\(\mathbb P_\beta\).  If the decoded value is a relative Random tag, a relative amoeba tag, a Cohen
tag, a direct coding tag \(\dot w\), or a well-order tag
\((\dot z_0,\dot z_1)\), then the next stage is the hybrid operation
prescribed in Definition~\ref{def:zero-local-allowable}, and
\(\dot I_{\beta+1}\) is the canonical lift of \(\dot I_\beta\).

It remains to describe the uniformization tag.  Suppose that
\[
   F(\beta)=(\dot x,\dot z,\dot m),
\]
where \(\dot x,\dot z\) are \(\mathbb P_\beta\)-names for reals and \(\dot m\)
is a \(\mathbb P_\beta\)-name for a natural number.  Let
\[
   A=\operatorname{supp}_{\mathrm{loc}}((\dot x,\dot m),\mathbb P_\beta).
\]
A candidate value in the test below is a \(\mathbb P_A\)-name \(\dot y_0\) for
a real, lifted canonically to a \(\mathbb P_\beta\)-name.  We distinguish two
alternatives.

\begin{enumerate}
\item[(a)] There are an ordinal \(\zeta<\rho\) and a \(\mathbb P_A\)-name
\(\dot y_0\) for a real such that
\[
   \mathbb P_\beta\Vdash
   (\dot x,\dot y_0^{\uparrow\beta})\in A_{\dot m},
\]
and no relative \(\zeta\)-allowable extension in \(W\) can force this pair
out of the relevant section.  More explicitly, for every relative
\(\zeta\)-allowable extension
\[
   (\mathcal S,\dot J)\triangleright_\zeta(\mathcal R_\beta,\dot I_\beta)
\]
in \(W\), write \(P_{\mathcal S}\) for the terminal post-branch forcing
associated with \(\mathcal S\).  If the displayed names are lifted to
\(P_{\mathcal S}\), then
\[
   \mathbb P_\beta\Vdash
   \text{``there is no condition in }
   P_{\mathcal S}/\dot G_\beta
   \text{ which forces }
   (\dot x^{\uparrow\mathcal S},\dot y_0^{\uparrow\mathcal S})
      \notin A_{\dot m^{\uparrow\mathcal S}}\text{.''}
\]
If such witnesses exist, choose first the least possible \(\zeta\), and then
choose the \(<_{M_1}\)-least localized presentation of a witnessing name
\(\dot y_0\).

In this case \(\dot y_0\) is declared to be a tentative value of rank
\(\zeta\).  The next forcing does not code this tentative value.  Instead it
codes another real on the same section: if the bookkeeping guess \(\dot z\) is
forced to be different from \(\dot y_0^{\uparrow\beta}\), use \(\dot z\);
otherwise use the \(<_{M_1}\)-least localized name \(\dot z'\) for a real
forced to be different from \(\dot y_0^{\uparrow\beta}\).  Denote the resulting
name by \(\dot u\).  The next coding operation is
\[
   \operatorname{Code}(\langle\dot x,\dot u,\dot m\rangle),
\]
understood in the hybrid sense of Definition~\ref{def:zero-local-allowable}.
The auxiliary name is updated by letting \(\dot I_{\beta+1}\) be the canonical
\(\mathbb P_{\beta+1}\)-name forced to be
\[
   \dot I_\beta^{\dot G_\beta}
   \cup
   \{(\dot x^{\dot G_\beta},
       (\dot y_0^{\uparrow\beta})^{\dot G_\beta},
       \dot m^{\dot G_\beta},\check\zeta)\}.
\]

\item[(b)] If clause~(a) fails, then the construction guesses.  The bookkeeping
guess \(\dot z\) is used, unless the bookkeeping code does not provide a usable
real name, in which case we use the \(<_{M_1}\)-least default real name.  The
next coding operation is
\[
   \operatorname{Code}(\langle\dot x,\dot z,\dot m\rangle),
\]
again understood as the corresponding hybrid product/iteration step, and no
new tentative value is added:
\[
   \dot I_{\beta+1}=\dot I_\beta
\]
up to the canonical lift to the longer hybrid presentation.
\end{enumerate}

If the decoded bookkeeping value is none of the recognized tags, the stage is
trivial and \(\dot I_{\beta+1}\) is the canonical lift of \(\dot I_\beta\).  If
there is a bookkeeping function \(F\in M_1\) such that \((\mathcal R,\dot I)\) is
\(\rho\)-allowable relative to \(F\), then we simply say that
\((\mathcal R,\dot I)\) is \(\rho\)-allowable.
\end{definition}

\begin{lemma}[Shrinking of allowability]
\label{lem:shrinking-allowability-local}
Work in \(W\).  Let \(\alpha<\beta\).  Let \(\mathcal R\) be a hybrid presentation,
and let \(\mathbb P\) be its terminal post-branch forcing.  Suppose that
\(\mathcal R\) is \(\beta\)-allowable, say there is a
\(\mathbb P\)-name \(\dot I^\beta\) such that \((\mathcal R,\dot I^\beta)\) is
\(\beta\)-allowable.  Then there is a \(\mathbb P\)-name \(\dot I^\alpha\) such
that \((\mathcal R,\dot I^\alpha)\) is \(\alpha\)-allowable.  In particular, the
classes of underlying allowable hybrid presentations are decreasing with the
rank of allowability.
\end{lemma}

\begin{lemma}[Allowability of tails]\label{lem:tail-allowability-local}
Let \((\mathcal R_\delta,\dot I_\delta)\) be \(\rho\)-allowable over \(W\),
witnessed by a bookkeeping function \(F\in M_1\), and let
\[
   (\mathcal R_\gamma,\dot I_\gamma),\qquad \gamma\leq\delta,
\]
be the canonical initial hybrid presentations and auxiliary names produced by
this construction.  If \(\beta\leq\gamma\leq\delta\), then
\[
   (\mathcal R_\gamma,\dot I_\gamma)
   \triangleright_\rho
   (\mathcal R_\beta,\dot I_\beta)
\]
as a relative \(\rho\)-allowable extension in \(W\).  Consequently, if
\(G_\beta\subseteq\mathbb P_\beta\) is generic over \(W\), then the quotient
\[
   \mathbb P_\gamma/G_\beta
\]
is a \(\rho\)-allowable tail over \(W[G_\beta]\).  If \(\xi<\rho\), then the
same quotient also has a \(W\)-presentation as a relative \(\xi\)-allowable
tail.  Finally, relative allowable extensions are transitive.
\end{lemma}

\begin{lemma}[Products of allowable local hybrid presentations]
\label{lem:rho-allowable-products-local}
Let \(\rho\) be an ordinal.  For \(i<2\), suppose that
\[
   (\mathcal R^i,\dot I^i)
\]
is \(\rho\)-allowable over \(W\), witnessed by a bookkeeping function
\(F_i\in M_1\).  Let \(H\subseteq Br\) be the fixed branch generic used to form
\(W=M_1[H]\).  For \(i<2\), let \(P^i\) be the terminal post-branch forcing over
\(W\) represented by \(\mathcal R^i\).  Then the ordinary product
\(P^0\times P^1\), computed over \(W\), has a canonical \(\rho\)-allowable local
hybrid presentation.

More precisely, there are a local hybrid presentation \(\mathcal R^\otimes\) over
\(W\), an auxiliary name \(\dot I^\otimes\) over its terminal post-branch forcing,
and a bookkeeping function \(F^\otimes\in M_1\), definable from \(F_0\) and
\(F_1\), such that
\[
   (\mathcal R^\otimes,\dot I^\otimes)
\]
is \(\rho\)-allowable over \(W\), witnessed by \(F^\otimes\), and its associated
terminal post-branch forcing is canonically forcing equivalent to
\(P^0\times P^1\).
\end{lemma}

\begin{proof}
We argue by induction on \(\rho\).  Write
\[
   \mathcal R^i=\langle \mathcal R^i_\beta\mid \beta\leq\delta_i\rangle
   \qquad (i<2),
\]
and let \(P^i_\beta\) denote the post-branch forcing over \(W\) represented by
\(\mathcal R^i_\beta\); thus \(P^i=P^i_{\delta_i}\).  The product presentation
\(\mathcal R^\otimes\) is obtained by first running \(\mathcal R^0\), and then
running a shifted copy of \(\mathcal R^1\).  Every name \(\dot\tau\) occurring in the second block is
replaced by its canonical lift \(\dot\tau^\uparrow\) to the corresponding
product initial segment, and all reservoir coordinates in the second block are
renamed to fresh coordinates.  Let \(P^\otimes_\gamma\) denote the post-branch forcing associated with the
initial product presentation \(\mathcal R^\otimes_\gamma\).  Thus, for every
\(\beta\leq\delta_1\),
\[
   P^\otimes_{\delta_0+\beta}
   \cong
   P^0\times P^1_\beta .
\]
The auxiliary names are defined by
\[
   \dot I^\otimes_{\delta_0+\beta}
   =
   \dot I^0\cup(\dot I^1_\beta)^\uparrow
\]
with the evident interpretation at stages \(\leq\delta_0\).  Limit stages use
the same mixed-support limits as the two original presentations.

For \(\rho=0\), this concatenation witnesses \(0\)-allowability.  The ordinary
Cohen, direct coding, well-order coding, and trivial stages are lifted
second-block stages.  If a shifted second-block stage is a relative regularity
stage with tag
\[
   (\mathrm{Rand};\dot a_0,\ldots,\dot a_{k-1}),
   \quad
   (\mathrm{Am}_{\mathcal N};\dot a_0,\ldots,\dot a_{k-1}),
   \quad\text{or}\quad
   (\mathrm{Am}_{\mathcal M};\dot a_0,\ldots,\dot a_{k-1}),
\]
then the product presentation uses the shifted tag obtained by lifting the
finitely many real names to the product initial segment.  If the original stage
has canonical admissibility witness
\[
   (A_{\mathrm{rb}},\dot W_{\mathrm{rb}},B^0_{\mathrm{rb}},\dot B_{\mathrm{rb}}),
\]
where
\[
   \dot W_{\mathrm{rb}}=L[T_2,\dot b_0,\ldots,\dot b_{k-1}]
\]
for the least tuple of names reading the displayed real parameters, then the
shifted stage has the lifted witness obtained from
\[
   A_{\mathrm{rb}}^\otimes,
   \qquad
   \langle \dot b_i^\uparrow\mid i<k\rangle,
   \qquad
   B^0_{\mathrm{rb}},
   \qquad
   \dot B_{\mathrm{rb}}^\uparrow.
\]
Here \(A_{\mathrm{rb}}^\otimes\) is the shifted copy of
\(A_{\mathrm{rb}}\) in the second block, each \(\dot b_i^\uparrow\) is the
canonical lift of \(\dot b_i\) to the product initial segment, and
\(\dot B_{\mathrm{rb}}^\uparrow\) is the canonical lift of the branch-footprint
name.  The lifted regularity base is therefore
\[
   L[T_2,\dot b_0^\uparrow,\ldots,\dot b_{k-1}^\uparrow]
\]
(with the evident meaning \(L[T_2]\) if \(k=0\)).  The Random or amoeba forcing
is computed in this lifted base, not in the larger product extension.  Hence the
displayed product equivalence is preserved at the shifted regularity stage.

Now assume \(\rho=\alpha+1\).  By the successor uniformization clause we mean
the two alternatives (a) and (b) in Definition~\ref{def:rho-allowable-local},
evaluated at a uniformization tag
\[
   (\mathrm{UF},\dot x,\dot z,\dot m).
\]
Thus clause~(a) searches through ranks \(\zeta<\rho\) for a localized name
\(\dot y\) which is forced into the relevant \(A_{\dot m}\)-section and cannot
be forced out by any relative \(\zeta\)-allowable tail; clause~(b) is used when
no such witness exists.

Consider a shifted uniformization stage \(\delta_0+\beta\), and let
\[
   (\mathrm{UF},\dot x,\dot z,\dot m)
\]
be the corresponding second-block bookkeeping value.  For \(\zeta<\rho\) and a
localized \(P^1_\beta\)-name \(\dot y\), write
\[
   \mathsf{Kick}^1_\beta(\zeta,\dot y)
\]
for the assertion that there are a relative \(\zeta\)-allowable extension
\[
   (\mathcal S,\dot J)\triangleright_\zeta
   (\mathcal R^1_\beta,\dot I^1_\beta)
\]
and, writing \(P_{\mathcal S}\) for the terminal post-branch forcing associated
with \(\mathcal S\), a quotient condition \(p\in P_{\mathcal S}/\dot G^1_\beta\)
such that
\[
   p\Vdash
   (\dot x^{\uparrow\mathcal S},
    \dot y^{\uparrow\mathcal S})
      \notin A_{\dot m^{\uparrow\mathcal S}}.
\]
Define \(\mathsf{Kick}^\otimes_\beta(\zeta,\dot y^\uparrow)\) analogously over
the product initial segment
\[
   (\mathcal R^\otimes_{\delta_0+\beta},
    \dot I^\otimes_{\delta_0+\beta}).
\]
Then
\[
   \mathsf{Kick}^1_\beta(\zeta,\dot y)
   \quad\Longleftrightarrow\quad
   \mathsf{Kick}^\otimes_\beta(\zeta,\dot y^\uparrow).
\]
The implication from left to right follows by taking the product of the fixed
first block with the displayed relative \(\zeta\)-allowable tail and applying
the induction hypothesis.  The implication from right to left follows by
restricting the witnessing product tail to the complete subpresentation
generated by the lifted second-block support of
\((\dot x,\dot y,\dot m)\), using Lemmas~\ref{lem:tail-allowability-local}
and~\ref{lem:shrinking-allowability-local}.

Similarly,
\[
   P^1_\beta\Vdash (\dot x,\dot y)\in A_{\dot m}
   \quad\Longleftrightarrow\quad
   P^0\times P^1_\beta
      \Vdash
      (\dot x^\uparrow,\dot y^\uparrow)\in A_{\dot m^\uparrow}.
\]
Hence the set of clause~(a) witnesses in the product construction is exactly
the lift of the set of clause~(a) witnesses in the second construction.  Thus
clause~(a) holds in the shifted product stage iff it holds in the original
second-block stage.  When it holds, the least rank and the \(<_{M_1}\)-least
localized witness are the lifted ones.  When it fails, clause~(b) uses the
lifted bookkeeping guess \(\dot z^\uparrow\).  The auxiliary name is updated by
\[
   \dot I^0\cup(\dot I^1_{\beta+1})^\uparrow.
\]
This proves the successor step.

Finally let \(\lambda\) be a nonzero limit ordinal.  If both inputs are
\(\lambda\)-allowable, then by the induction hypothesis the product
presentation is \(\zeta\)-allowable for every \(\zeta<\lambda\).  Since
\[
   \mathcal A^W_\lambda
   =
   \bigcap_{\zeta<\lambda}\mathcal A^W_\zeta,
\]
the product presentation is \(\lambda\)-allowable.  The bookkeeping witness at
a limit level is obtained by packaging the lower-level witnesses together with
the concatenation description.  Thus \(F^\otimes\in M_1\) is definable from
\(F_0\) and \(F_1\).
\end{proof}

\end{document}
